% Vista científica exacta materializada desde una rama condicional.
% Fuente: source/demostraciones_primarias/09_06h_atlas_familias_rev6.tex; rama: HMTIncludeAtlasStructural.
% No modifica el contenido matemático de la rama de origen.
\chapter[Classification nonadique des familles]{Classification nonadique des familles de trajectoires}
\label{chap:atlas-genealogia-masas}
\index{atlas nonadique}
\index{famille de chemin}
\index{multisection}
\index{secteur nul}

L'atlas ne commence pas par une liste de particules ni par une table de masses.
Il commence par un support fini orienté et par les lecteurs qui conservent
l'histoire de ses chemins. Sa fonction est de répondre à une question antérieure à
toute comparaison métrologique : quelles classes de persistance la géométrie nonadique peut-elle
soutenir et quelle information doit être conservée pour les distinguer ?

La réponse comporte trois cardinaux, mais non trois inventaires rivaux :
\[
 81\ \text{cellules},\qquad
 56\ \text{familles de routes},\qquad
 13\ \text{multisections famille--niveau}.
\]
Le premier nombre compte des positions orientées ; le deuxième identifie les
positions qui publient la même signature finie de chemin ; le troisième réunit les
positions qui partagent famille et profondeur par rapport au centre. Aucun
d'eux ne compte des masses connues, des noms d'espèces ou des lignes d'une
comparaison externe.

\section{La géométrie complète des \(81\) cellules}

Soit
\[
 I_9=\{1,\ldots,9\},
 \qquad
 \mathcal C_{81}=I_9\times I_9.
\]
Les coordonnées \((r,c)\) conservent les deux directions d'APP\@. La parité
de chaque coordonnée définit quatre secteurs internes :
\[
 \operatorname{fam}(r,c)=
 \begin{cases}
  \ell^- ,& r\equiv1,\ c\equiv1\pmod2,\\
  \nu ,& r\equiv1,\ c\equiv0\pmod2,\\
  d ,& r\equiv0,\ c\equiv1\pmod2,\\
  u ,& r\equiv0,\ c\equiv0\pmod2.
 \end{cases}
 \tag{A1}\label{eq:familias-paridad-atlas}
\]
La profondeur dans l'atlas n'est pas ajoutée comme étiquette indépendante. C'est la
distance de Tchebychev au centre :
\[
 G(r,c)=\max\{|r-5|,|c-5|\}.
 \tag{A2}\label{eq:generacion-chebyshev-atlas}
\]
Par conséquent, famille et niveau proviennent de la position elle-même. En particulier,
le centre n'est pas obtenu en sélectionnant une espèce de l'extérieur : c'est l'unique
cellule de profondeur zéro.

\begin{theorem}[Recensement des familles et des générations de l'atlas]
\label{thm:censo-atlas-81}
Les applications \eqref{eq:familias-paridad-atlas} et
\eqref{eq:generacion-chebyshev-atlas} décomposent l'atlas de manière
disjointe comme
\[
 81=45_{\rm leptones}+36_{\rm quarks}
   =25_{\ell^-}+20_{\nu}+20_{d}+16_{u}.
 \tag{A3}\label{eq:censo-81-cuatro-familias}
\]
En raffinant par \(G\), on obtient exactement les treize multisections
\[
\begin{array}{c|c|c}
\text{famille}&\text{niveaux}&\text{cardinaux}\\ \hline
\ell^-&G0,G2,G4&1,8,16\\
\nu&G1,G2,G3,G4&2,4,6,8\\
d&G1,G2,G3,G4&2,4,6,8\\
u&G1,G3&4,12.
\end{array}
\tag{A4}\label{eq:trece-multisecciones-atlas}
\]
La cellule \((5,5)\) est l'unique composante de
\(\mathfrak S_{\ell^-,0}\).
\end{theorem}

\begin{proof}
Dans \(I_9\), il y a cinq coordonnées impaires et quatre paires. Les quatre classes
de parité ont donc pour cardinaux
\[
 5\cdot5=25,\qquad
 5\cdot4=20,\qquad
 4\cdot5=20,\qquad
 4\cdot4=16,
\]
ce qui prouve \eqref{eq:censo-81-cuatro-familias} ; les deux premières
classes réunissent \(45\) cellules et les deux dernières, \(36\).

Pour le raffinement, nous faisons croître des carrés centrés en \((5,5)\).
Dans la classe impair--impair, il y a une cellule au rayon zéro, \(3^2-1=8\)
dans la couche de rayon deux et \(5^2-3^2=16\) dans celle de rayon quatre.
Dans la classe impair--pair, les rectangles accumulés jusqu'aux rayons
\(1,2,3,4\) contiennent, respectivement,
\[
 1\cdot2=2,\qquad 3\cdot2=6,\qquad
 3\cdot4=12,\qquad5\cdot4=20
\]
cellules ; leurs différences successives sont \(2,4,6,8\). La classe
pair--impair donne le même calcul avec les coordonnées échangées. Dans la
classe pair--pair apparaissent \(2^2=4\) cellules au rayon un et
\(4^2-2^2=12\) au rayon trois. Ces couches sont disjointes et épuisent
\(I_9^2\), de sorte qu'il ne reste aucune autre multisection.
\end{proof}

Le mot \emph{niveau} désigne dans ce chapitre la profondeur combinatoire
\(G\). La lecture physique ultérieure conserve cette géométrie, mais ne la
définit pas. On évite ainsi de transformer les noms des particules en prémisses
du recensement qui doit les expliquer.

\section{Inversion centrale, \(41\) orbites et unicité du centre}

L'inversion de l'atlas est
\[
 \rho_{\rm c}(r,c)=(10-r,10-c).
 \tag{A5}\label{eq:inversion-central-atlas-rev6}
\]
Comme \(10-r\) a la même parité que \(r\), l'inversion conserve les
quatre familles. En outre,
\[
 |(10-r)-5|=|r-5|,
 \qquad
 |(10-c)-5|=|c-5|,
\]
et, par conséquent, conserve le niveau \(G\). Chacune des treize
multisections est \(\rho_{\rm c}\)-stable.

\begin{proposition}[Orbites et exception centrale]
\label{prop:orbitas-centro-atlas-rev6}
L'action de \(\rho_{\rm c}\) sur \(\mathcal C_{81}\) possède exactement
\(41\) orbites. Son unique point fixe est \((5,5)\). Par conséquent, la
multisection \(\mathfrak S_{\ell^-,0}\) admet l'unique sélecteur ponctuel
\(\rho_{\rm c}\)-équivariant ; dans les douze autres, l'objet naturel est
l'orbite ou la multisection complète.
\end{proposition}

\begin{proof}
L'équation
\[
 (r,c)=\rho_{\rm c}(r,c)
\]
équivaut à \(2r=2c=10\), donc \(r=c=5\). Les quatre-vingts autres cellules se
regroupent en quarante couples, d'où \(1+40=41\) orbites.

Soit \(y\) une multisection non centrale et supposons qu'un sélecteur ponctuel
équivariant choisisse \(a(y)\in\mathfrak S_y\). Comme \(\rho_{\rm c}\) agit
trivialement sur le type famille--niveau, l'équivariance imposerait
\[
 a(y)=a(\rho_{\rm c}y)=\rho_{\rm c}a(y).
\]
Le point choisi devrait être fixe, mais l'unique point fixe est dans
\(\mathfrak S_{\ell^-,0}\). La contradiction démontre l'affirmation.
\end{proof}

Cette proposition explique pourquoi une espèce interne n'est pas, en général, une
case isolée. L'électron central est exceptionnel parce que sa multisection
se réduit à un point. Dans les autres familles, l'orientation conjuguée fait
partie de l'identité qui doit être conservée.

\section{Les \(56\) familles de chemins}

Chaque cellule porte l'ombre finie de son histoire de transport :
\[
 \mathcal R(r,c)=
 \bigl(
 n_A,\ s_C,\ k_{\Delta_4},\
 \nu_{120},\ \nu_{270},\ \tau_{12}
 \bigr)_{r,c}.
 \tag{A6}\label{eq:lector-familia-ruta-atlas}
\]
Les cinq premières coordonnées sont entières ; \(\tau_{12}\) est le mot
dodécaphasique orienté. Deux cellules appartiennent à la même famille de chemins
si et seulement si elles ont le même sextuplet \(\mathcal R\).

\begin{theorem}[Recensement des familles de chemins et mémoire minimale]
\label{thm:censo-56-familias-ruta-rev6}
L'image de \(\mathcal R\) contient exactement \(56\) éléments.
L'histogramme des tailles de ses fibres est
\[
 41\times1,\qquad
 10\times2,\qquad
 2\times3,\qquad
 1\times4,\qquad
 2\times5.
 \tag{A7}\label{eq:histograma-fibras-ruta}
\]
En particulier,
\[
 41+10+2+1+2=56,
\]
\[
 41\cdot1+10\cdot2+2\cdot3+1\cdot4+2\cdot5=81.
 \tag{A8}\label{eq:doble-censo-rutas}
\]
L'ordre lexicographique étant fixé dans chaque fibre, le rang
\[
 \kappa(r,c)\in\{0,1,2,3,4\}
\]
rend injective l'application
\[
 (r,c)\longmapsto\bigl(\mathcal R(r,c),\kappa(r,c)\bigr).
 \tag{A9}\label{eq:lift-memoria-cinco-atlas}
\]
Cinq est la taille minimale d'un alphabet uniforme capable d'effectuer ce
relèvement.
\end{theorem}

\begin{proof}
L'énumération exhaustive de \(I_9^2\) regroupe les quatre-vingt-un sextuplets
\eqref{eq:lector-familia-ruta-atlas} par égalité coordonnée par coordonnée.
Elle produit les cinq tailles de fibre de
\eqref{eq:histograma-fibras-ruta} ; les deux identités
\eqref{eq:doble-censo-rutas} vérifient simultanément le nombre de classes
et la couverture des quatre-vingt-une cellules.

Le rang lexicographique distingue les éléments de chaque fibre, de sorte que
\eqref{eq:lift-memoria-cinco-atlas} est injective. Comme il existe deux fibres
de cardinal cinq, tout alphabet commun comportant moins de cinq symboles
identifierait deux cellules de l'une d'elles par le principe des tiroirs.
\end{proof}

Le rang \(\kappa\) est une mémoire de position au sein d'une fibre ; il n'est
ni charge, ni masse, ni couleur. Sa fonction est de retrouver la cellule que la signature visible
du chemin avait identifiée à d'autres. Le nombre \(56\) compte les classes du
sextuplet \(\mathcal R\), tandis que le nombre \(13\) compte les classes du couple
\((\operatorname{fam},G)\). Ce sont des quotients différents du même atlas.

\paragraph{Quotients cellulaire et orbital de l'atlas : domaines et applications de projection.}
La projection brute CT--108 et la classification scientifique lisent des données
distinctes des mêmes quatre-vingt-une cellules. Si
\(q_{\rm CT}\) conserve seulement la forme du mot tritique
dodécaphasique, tandis que \(\mathcal R\) conserve le sextuplet de chemin et
\(q=(\operatorname{fam},G)\) conserve famille et profondeur, le tableau
typé est
\[
\begin{array}{rcll}
 \mathcal C_{81}
 &\xrightarrow{\ q_{\rm CT}\ }&
 \mathcal W^{\rm raw}_{7}
 & \text{formes brutes CT--108},\\[2mm]
 \mathcal C_{81}
 &\xrightarrow{\ \mathcal R\ }&
 \Sigma^{\rm ruta}_{56}
 & \text{familles de routes},\\[2mm]
 \mathcal C_{81}
 &\xrightarrow{\ q\ }&
 \Sigma_{13}
 & \text{multisections famille--niveau}.
\end{array}
\tag{A9a}\label{eq:dos-cocientes-atlas-rev8}
\]
Les sept fibres brutes ont pour multiplicités
\[
 54,2,2,2,7,7,7,
 \qquad54+3\cdot2+3\cdot7=81.
\]
Par conséquent, le recensement archéologique \(81\to7\) ne remplace pas l'atlas
\(81/56/13\). L'écriture abrégée \(81/56/13\) n'affirme pas non plus qu'il existe
une flèche canonique \(56\to13\) : certaines fibres de \(\mathcal R\) traversent
plus d'un couple famille--niveau. Les applications \(\mathcal R\) et \(q\) sont des
quotients parallèles et conservent des invariants différents. Dans le secteur quark,
la convention active reste
\[
 \operatorname{col}(r,c)=r-c\pmod3;
\]
l'écriture historique \(c-r\) est sa convention conjuguée, non une autre couleur ni
un autre quotient de l'atlas.

\section{Distinction entre cellule, famille de chemins et multisection}

Définissons
\[
 q:\mathcal C_{81}\longrightarrow\Sigma_{13},
 \qquad
 q(r,c)=\bigl(\operatorname{fam}(r,c),G(r,c)\bigr),
 \tag{A10}\label{eq:cociente-trece-multisecciones}
\]
et, pour \(y\in\Sigma_{13}\),
\[
 \mathfrak S_y^{\rm atlas}
 =q^{-1}(y)
 =\{(r,c)\in\mathcal C_{81}:q(r,c)=y\}.
 \tag{A11}\label{eq:multiseccion-atlas-rev6}
\]
Les quatre niveaux de lecture sont alors séparés :

\begin{enumerate}[label=(\roman*)]
 \item une \emph{cellule} est une position orientée \((r,c)\), avec toutes ses
 données locales ;
 \item une \emph{famille de chemins} est une fibre de \(\mathcal R\) : elle conserve
 une même signature finie de transport ;
 \item une \emph{multisection} est une fibre de \(q\) : elle conserve famille,
 profondeur et totalité de ses orientations conjuguées ;
 \item une \emph{représentation physique} est la réalisation ultérieure d'une
 multisection enrichie dans un espace d'états ; ce n'est pas une autre partition du
 tableau ni le choix arbitraire de l'une de ses cellules.
\end{enumerate}

Cette distinction donne un contenu mathématique au mot \emph{espèce}. Le type
interne est la multisection. Une particule concrète est une section
persistante réalisée en elle. La représentation physique conserve l'action,
l'orientation et les observables du type ; elle ne revient pas en arrière pour redéfinir
l'atlas à partir d'une masse déjà connue.

\section{Couleur résiduelle comme orientation ternaire}

Dans le secteur quark, formé des \(36\) cellules avec \(r\) pair, nous définissons
\[
 \operatorname{col}(r,c)=r-c\pmod3,
 \tag{A12}\label{eq:color-residual-atlas-rev6}
\]
avec le dictionnaire
\[
 0\longmapsto R,\qquad1\longmapsto G,\qquad2\longmapsto B.
\]

\begin{proposition}[Bilan de couleur résiduelle]
\label{prop:color-atlas}
Les trente-six cellules de quark se répartissent exactement comme
\[
 36=12_R+12_G+12_B.
 \tag{A13}\label{eq:balance-color-atlas-rev6}
\]
L'inversion centrale fixe \(R\) et échange \(G\) et \(B\).
\end{proposition}

\begin{proof}
L'une quelconque des quatre valeurs paires de \(r\) étant fixée, lorsqu'on parcourt
\(c=1,\ldots,9\), chaque classe modulo trois apparaît trois fois. Chaque résidu
a donc \(4\cdot3=12\) représentants. En outre,
\[
 \operatorname{col}\bigl(\rho_{\rm c}(r,c)\bigr)
 =(10-r)-(10-c)
 =-(r-c)\pmod3.
\]
La négation modulaire fixe la classe zéro et échange les classes un et deux.
\end{proof}

Le tritique \(R,G,B\) est donc un observable exact de l'atlas et non une
décoration ajoutée à une table de quarks. La représentation de couleur se
construit ensuite sur ce support ternaire ; le recensement \(12+12+12\) est sa
géométrie préalable.

\section{De l'extension survivante à la signature}

L'atlas classe des coupes finies. L'identité d'une particule requiert
l'histoire profonde dont cette coupe provient. Soit
\[
 \Omega_{\rm HMT}^{\rm surv}
 =\varprojlim_m\mathcal S_m
\]
le système inverse des états survivants. La projection sur l'atlas \(\pi_{81}\) et l'ombre
du chemin forment la chaîne
\[
 \Omega_{\rm HMT}^{\rm surv}
 \xrightarrow{\ \operatorname{sh}_{108}\ }
 \Gamma_{108}^{\rm enr}
 \xrightarrow{\ \mathcal E\ }
 \mathcal L_{108}
 \xrightarrow{\ L_{\rm tick}\ }
 \mathbb Z^5.
 \tag{A14}\label{eq:cadena-extension-firma-atlas-rev6}
\]
Ici, \(\operatorname{sh}_{108}\) publie un chemin de \(108\) pas sans
oublier feuillet, orientation, frontière ni prédécesseurs ; \(\mathcal E\) forme le
registre dodécaphasique ; et \(L_{\rm tick}\) extrait
\[
 v(x)=
 \bigl(n_A,s_C,k_{\Delta_4},\nu_{120},\nu_{270}\bigr).
 \tag{A15}\label{eq:firma-cinco-extension-rev6}
\]
Aucune de ces trois flèches ne consulte une masse, une unité ou le nom
d'une particule.

La raison de conserver le registre avant la projection se voit dans ses douze
événements orientés \(e_0,\ldots,e_{11}\). Pour \(0\leq i\leq5\), soit
\[
 p_i=e_i+e_{i+6},
 \qquad
 a_i=e_i-e_{i+6},
 \qquad
 s_C=\sum_{i=0}^{5}p_i,
\]
et, pour \(0\leq i<5\),
\[
 M_i=p_i-p_5.
\]
Alors
\[
 p_5=\frac{s_C-\sum_{i=0}^{4}M_i}{6},
 \qquad
 p_i=p_5+M_i,
\]
\[
 e_i=\frac{p_i+a_i}{2},
 \qquad
 e_{i+6}=\frac{p_i-a_i}{2}.
 \tag{A16}\label{eq:reconstruccion-registro-doce}
\]
L’identité
\[
 1+5+6=12
\]
est ainsi une reconstruction réversible : un total, cinq différences et six
orientations retrouvent les douze événements. Projeter sur la signature de cinq
entiers avant de composer les histoires éliminerait précisément le placement et
l'orientation qui permettent de décider si deux chemins peuvent être composés.

Le quotient \(q\) se relève à l'histoire profonde au moyen de la projection sur l'atlas
\(\pi_{81}:\Omega_{\rm HMT}^{\rm surv}\to\mathcal C_{81}\) :
\[
 \widetilde{\mathfrak S}_y
 =\{x\in\Omega_{\rm HMT}^{\rm surv}:q(\pi_{81}x)=y\}.
 \tag{A17}\label{eq:multiseccion-enriquecida-rev6}
\]
C'est la multisection enrichie. Son ombre appartient à l'atlas ; son chemin
appartient à \(\Gamma_{108}^{\rm enr}\) ; sa mémoire complète appartient au
registre ; et son évaluation n'est appliquée qu'à la fin. La généalogie n'est pas une
note ajoutée à la valeur : c'est l'objet qui rend la valeur possible et conserve les
autres lectures de la même section.

\section{Bifurcation du secteur nul avant le caractère positif}

Le caractère de masse agit sur la signature entière et a un codomaine positif :
\[
 \chi_{\rm mass}:\mathbb Z^5\longrightarrow\mathbb R_{>0}.
 \tag{A18}\label{eq:caracter-positivo-atlas-rev6}
\]
Par définition, il n'existe pas de signature \(v\) avec
\(\chi_{\rm mass}(v)=0\). Le secteur nul n'est donc pas une
quatorzième multisection ni la limite dégénérée des treize précédentes.

Son support algébrique appartient à la branche scindée
\[
 \mathcal A_{\rm sp}
 =\mathbb R[\mathsf R]/(\mathsf R^2-1),
 \qquad
 P_\pm=\frac{1\pm\mathsf R}{2}.
\]
Tout élément s'écrit de manière unique comme
\[
 z=r_+P_++r_-P_-,
 \qquad
 N_{\rm sp}(z)=r_+r_-,
\]
et les deux droites
\[
 \mathbb R P_+,\qquad\mathbb R P_-
\]
sont nulles parce que \(N_{\rm sp}(\lambda P_\pm)=0\). Une section nulle est un
chemin ouvert transporté dans l'une de ces directions. L'intérieur matériel
requiert le couplage des deux, pour lequel \(r_+r_->0\).

Ainsi sont typées deux branches :
\[
 \text{persistance matérielle}
 \longrightarrow\mathbb Z^5
 \xrightarrow{\chi_{\rm mass}}\mathbb R_{>0},
\]
\[
 \text{propagation nulle}
 \longrightarrow\partial_{\rm null}\mathcal C_{\rm sp}^{+}.
\]
Les réalisations photonique et de jauge lisent ensuite la seconde branche. Sa
séparation du caractère positif empêche d'effacer photon et gluon parce qu'ils n'apparaissent pas
dans une liste de masses, et évite en même temps de fabriquer une masse nulle au moyen d'un
caractère dont le codomaine l'exclut.

\section{Système fini de classification et mémoire des trajectoires}

Les trois quotients du chapitre remplissent des fonctions complémentaires :
\[
 \mathcal C_{81}
 \xrightarrow{\ \mathcal R\ }\mathcal R_{56},
 \qquad
 \mathcal C_{81}
 \xrightarrow{\ q\ }\Sigma_{13},
 \qquad
 \mathcal C_{81}
 \xrightarrow{\ /\langle\rho_{\rm c}\rangle\ }
 \mathcal O_{41}.
\]
Le premier conserve la signature du chemin ; le deuxième conserve famille et niveau ;
le troisième conserve la conjugaison centrale. Avec la couleur résiduelle et le
relèvement \(\kappa\), ils forment un système fini capable de reconstruire
quelle information est visible, laquelle est identifiée et laquelle doit subsister comme
mémoire.

La chaîne causale de l'atlas est donc fixée :
\[
 \boxed{
 \begin{gathered}
 \text{extension survivante}
 \longrightarrow\text{route enrichie}
 \longrightarrow\text{registre}
 \longrightarrow\text{signature}\\
 \longrightarrow\text{caractère}
 \longrightarrow\text{lecture dimensionnelle}
 \end{gathered}}
 \tag{A19}\label{eq:cadena-rectora-atlas-rev6}
\]
Les quatre-vingt-une cellules constituent le domaine local ; les cinquante-six
familles expriment des coïncidences de chemins ; les treize multisections sont les
types internes ; les quarante-et-une orbites conservent la conjugaison ; et le
centre \((5,5)\) est l'unique exception ponctuelle. La masse apparaît ensuite
comme lecture d'une persistance déjà construite. L'atlas n'est pas une table
abrégée de résultats : c'est la structure qui explique pourquoi une table peut
exister.
